% Generated draft: requires mathematical and source review.
\documentclass[11pt]{article}
\usepackage[margin=25mm]{geometry}
\usepackage{fontspec}
\setmainfont{Latin Modern Roman}
\usepackage{amsmath,amssymb,mathtools,unicode-math}
\setmathfont{Latin Modern Math}
\usepackage{enumitem,needspace,xurl,hyperref}
\setlength{\parindent}{0pt}
\setlength{\parskip}{5pt}
\setlength{\emergencystretch}{3em}
\newcommand{\N}{\mathbb{N}}
\newcommand{\Z}{\mathbb{Z}}
\newcommand{\Q}{\mathbb{Q}}
\newcommand{\R}{\mathbb{R}}
\newcommand{\C}{\mathbb{C}}
\newcommand{\F}{\mathbb{F}}
\newcommand{\A}{\mathbb{A}}
\newcommand{\PP}{\mathbb{P}}
\newcommand{\E}{\mathbb{E}}
\newcommand{\eps}{\varepsilon}
\newcommand{\dd}{\,\mathrm{d}}
\begin{document}
\section*{UnsolvedMath ID: 701694}
\textbf{Source problem code:} catalogue entry
\subsection{Definitions and mathematical statement}
Let \(F_n=\langle x_1,\ldots,x_n\rangle\) be the free group of rank \(n\), for each positive integer \(n\). For a subset \(S\) of a group \(F\), write \(\langle\!\langle S\rangle\!\rangle_F\) for its normal closure in \(F\). Define
\[
N_n=\langle\!\langle w^4:w\in F_n\rangle\!\rangle_{F_n},
\qquad B(n,4)=F_n/N_n.
\]
Thus \(B(n,4)\) is the free Burnside group on \(n\) generators subject to the identity \(g^4=1\) for every element \(g\).

For \(F=F_3\), put \(N=N_3\) and define
\[
\mu=\min\left\{r\in\Z_{\geq0}:\ \exists w_1,\ldots,w_r\in F,\quad
N=\langle\!\langle w_1^4,\ldots,w_r^4\rangle\!\rangle_F\right\},
\]
with \(\mu=\infty\) if this set is empty. The equality of normal subgroups means that
\[
\langle x_1,x_2,x_3\mid w_1^4=1,\ldots,w_r^4=1\rangle
\longrightarrow B(3,4)
\]
is an isomorphism under the natural map sending each displayed generator to its canonical image. Each fourth-power relator is counted once in the displayed list, irrespective of its word length.

\textbf{Questions (catalogue entry).}
\begin{enumerate}
\item[(a)] Is \(\mu<10^6\)?
\item[(b)] Is \(\mu<10^3\)?
\item[(c)] What is the exact value of \(\mu\)?
\end{enumerate}

The supplied source reports that \(|B(3,4)|=2^{69}\), attributing this to Bayes, Kautsky, and Wamsley (1974), whose proof uses Sanov's (1940) finiteness theorem for \(B(n,4)\). It further reports that Sanov's proof for \(B(3,4)\) uses more than \(2^{32}\) fourth powers, explaining this by the fact that the subgroup generated by two canonical generators and the square of the third has order \(2^{32}\). Finally, it attributes to Havas and Newman (1980) the lower bound of \(105\) defining relations for \(B(3,4)\); in particular, the reported bound gives \(\mu\geq105\).

The formulation above interprets ``defined with fourth powers'' as a presentation on the three free generators with relators that are fourth powers of words. The source does not explicitly specify the presentation conventions; this interpretation requires human confirmation. All historical claims and the open status are supplied source assertions, not independently verified conclusions.
\subsection{Short English statement}
For the free Burnside group \(B(3,4)\), determine the smallest number of fourth-power relations needed on its three free generators. In particular, can fewer than \(1{,}000{,}000\), or fewer than \(1{,}000\), such relations suffice? The supplied record labels these questions open and reports a lower bound of \(105\) defining relations.
\subsection{Sources}
\begin{itemize}
\item Local UnsolvedMath registry, permanent ID \(701694\), exact source code \texttt{catalogue entry}; supplied complete source record, upstream revision \texttt{d371ee917ac9647fbb06cd6fc529deec6f3dda82}.
\item \emph{The Kourovka Notebook (Unsolved Problems in Group Theory)}, Problem 20.36; proposed by G. Havas and M. F. Newman (2022). The supplied background identifies the 20th issue (2022) and the 21st issue edition (2026), with identifier \texttt{arXiv:1401.0300v48}. The supplied collection description names E. I. Khukhro and V. D. Mazurov as editors.
\item Bayes, Kautsky, and Wamsley (1974): cited in the supplied statement for the order \(2^{69}\) of \(B(3,4)\). No title, publication details, or link is supplied.
\item Sanov (1940): cited in the supplied statement for finiteness of \(B(n,4)\), with the additional reported qualification concerning more than \(2^{32}\) fourth powers in the proof for \(B(3,4)\). No title, publication details, or link is supplied.
\item Havas and Newman (1980): cited in the supplied statement for the lower bound of \(105\) defining relations. No title, publication details, or link is supplied.
\end{itemize}
\end{document}
